\documentclass[a4paper,12pt]{report}

\usepackage[utf8]{inputenc}
\usepackage{CJKutf8}
\usepackage[T1]{fontenc}
\usepackage{amsmath,amssymb}
\usepackage{graphicx}
\usepackage{geometry}
\geometry{margin=25mm}

\begin{document}
\begin{CJK}{UTF8}{min}

\chapter*{第1章　序論}
\addcontentsline{toc}{chapter}{第1章　序論}

近年、人工知能（Artificial Intelligence: AI）および機械学習は、金融、流通、交通など、時間とともに状況が変化する分野の予測に利用されている。とりわけ株式市場では、過去の価格、取引量、企業情報および経済指標を組み合わせ、将来の価格変動や売買時期を推定する研究が発展してきた。これに対し、コンサートチケットの二次流通市場では、ウェブ上に多数の出品情報が存在するにもかかわらず、それらを継続的に収集し、価格や需要を定量的に予測する仕組みは十分に整備されていない。購入者は、現在の出品価格が高いか安いか、公演日まで待てば価格が下がるか、同等条件の安価な出品が新たに現れるかを、過去の経験、SNS上の情報および個人の感覚から判断する場合が多い。

チケットの価格は、出演者の人気だけで決まるものではない。公演日、曜日、開演時間、会場、会場収容人数、ツアー内の公演位置、公演日までの残り日数、出品枚数、チケット種別、名義種別、受渡方法など、複数の条件が価格形成に関係すると考えられる。さらに、STARTO ENTERTAINMENT所属グループを中心とする公演の二次流通では、出品説明文に記載される番手、名義数、ランダム配布、制作開放席、見切れ・注釈付き席、本人確認、QRコード単位の取引条件なども、同一公演内の価格差を生じさせる可能性がある。これらの情報には数値、カテゴリ、日時および自由記述文が混在するため、人手だけで多数の出品を同じ基準から比較することは難しい。

そこで本研究では、チケット二次流通プラットフォーム「チケテン」に掲載された公開情報を反復的に収集し、チケットごとの成立価格、需要状態および購入時期を検討するためのデータセットとAIモデルを構築する。対象は、二次流通市場で継続的な出品が確認できるSTARTO ENTERTAINMENT所属グループを中心とする複数公演である。チケテンからは、価格、公演日時、会場、枚数、出品状態、出品説明文等を取得する。また、チケテンだけでは把握できないファンクラブ会員数、チケット定価、会場収容人数、ツアー公演数、ツアー初日および最終公演等を手動データとして整備し、出品データと関連付ける。これにより、一件のチケットを価格、商品条件、時間、需要と供給という複数の観点から表現する。

本研究では、第一に、売却済みチケットから観測される価格を用いて、類似条件に基づく妥当価格を推定する。第二に、各観測時点から一定期間内に出品が継続するか、売却済みとなるか、削除されるかを区別して需要状態を推定する。第三に、同一公演・同等条件で現在価格より安い新規出品が現れる可能性を推定し、これらを統合して「購入する」または「待つ」という判断を支援する。ただし、本研究でいう妥当価格は、チケットの真正な価値や取引の正当性を決定する価格ではなく、収集した市場データと設定した条件の範囲で得られる統計的な参考値である。また、削除された出品を売却済みとみなさず、予測時点より後に判明する情報を学習時の説明変数へ混入させない。本研究は、経験や勘に依存しやすい購入判断を、再現可能なデータ分析によって補助するための基礎を構築するものである。

\chapter*{第2章　研究背景}
\addcontentsline{toc}{chapter}{第2章　研究背景}

チケット市場は、主催者や公式販売事業者が販売する一次流通市場と、購入されたチケットが別の利用者へ譲渡される二次流通市場に分けられる。二次流通では、出品者が提示する価格と購入者の需要が取引時点ごとに変化するため、同じ出演者、同じ会場、同じ公演日のチケットであっても価格は一定ではない。人気公演では定価を上回る価格で出品される一方、公演日が近づいた場合、売り切るための値下げや安価な新規出品が生じることもある。しかし、チケットには公演日時という明確な期限があり、公演終了後はその公演へ入場する権利としての価値を失う。この点で、継続的に取引される株式等の金融商品とは異なり、購入を待つことで安く買える可能性と、希望条件のチケットを買い逃す危険が同時に存在する。

株式市場の価格予測では、過去の価格系列、出来高および外部情報を説明変数とし、時系列分析、回帰分析、決定木およびニューラルネットワーク等を用いて将来の値を推定する。チケット二次流通市場にも、価格が需要と供給に応じて変化し、過去の取引や市場状況が次の判断材料になるという共通点がある。一方、チケット市場では、取引対象が公演ごとに異なり、販売期間が限られ、出品条件が自由記述文に含まれる。加えて、表示から消えた出品が売れたのか、出品者に取り下げられたのかを区別できない場合がある。そのため、金融市場の予測方法をそのまま適用するのではなく、公演期限、商品条件および観測可能な出品状態を考慮したデータ設計が必要である。

チケテンの出品情報には、価格、公演日時、会場、枚数、チケット種別、名義種別、受渡方法および出品状態などの表形式情報に加え、出品者が条件を記載した説明文が存在する。特に本研究の対象市場では、「複数名義中の何番手か」「QRコードごとの配布か」「ランダム配布か」「制作開放席か」「見切れ・注釈付き席か」など、価格へ影響し得る情報が説明文にのみ記載される場合がある。単純なキーワードの有無だけでは、数値と条件の対応や否定表現を正確に扱えない可能性がある。一方、会場収容人数、ツアー全体の公演数、チケット定価およびファンクラブ会員数などは、個々の出品ページだけからは得られない。したがって、ウェブから自動収集する構造化データと説明文に、公開情報を基に手動で整備した公演・会場・アーティスト情報を組み合わせる必要がある。

本研究で用いる\texttt{tiketen\_scraper}は、チケテンから出品中、売却済みおよび削除されたチケットの状態を反復的に取得し、グループ別のマスターデータと月別・公演別のスナップショットを作成する。継続観測によって、単一時点の価格表だけではなく、出品が時間の経過とともにどの状態へ移ったか、市場内の出品数や価格分布がどのように変化したかを扱える。ただし、出品情報の識別子変更、初回取得以前に売却済みであったチケットの売却時刻、欠測期間中の状態変化などは、誤った教師データを生む原因となる。このため、同一出品の論理的な統合、状態変化を実際に観測できたデータの識別、欠測や異常スナップショットの検査が不可欠である。

以上の背景から、チケット二次流通の分析には、価格回帰だけでなく、出品状態の確率推定と、将来の代替出品の出現予測を分けて扱う必要がある。価格モデルは「類似条件ならいくらで取引されるか」を示すが、それだけでは現在買うべきかを決められない。現在の出品が短期間に売れる可能性、より安い代替が現れる可能性、および待つことによる機会損失を合わせて評価して初めて、購入時期に関する判断材料となる。本研究は、期限付きで条件の異質性が大きいチケット市場に対し、構造化情報、説明文、手動データおよび時系列の出品状態を統合する点に特徴がある。

\chapter*{第3章　研究目的}
\addcontentsline{toc}{chapter}{第3章　研究目的}

本研究の目的は、STARTO ENTERTAINMENT所属グループを中心とする公演のチケット二次流通を対象として、チケテンから継続収集した出品データと手動で整備した補足データを統合し、チケットの妥当価格、需要状態および購入時期を推定するAIモデルの基礎を構築することである。ここで妥当価格とは、市場で成立し得る唯一の正解価格を意味するものではない。本研究で収集した期間、対象グループ、公演および取引条件の範囲において、売却済みチケットと類似する条件から推定される一枚当たりの参考価格と定義する。また、買い時とは価格が将来必ず最小となる時点ではなく、現在のチケットを買い逃す損失と、待つことで安価な選択肢を得る可能性を比較した意思決定上の判断とする。

第一の目的は、価格形成に関係する複数種類の情報を統合し、売却価格を予測することである。チケテンから取得する価格、公演日時、会場、枚数、チケット種別、名義種別、受渡方法、出品者情報および説明文に加え、手動で収集する定価、会場収容人数、ファンクラブ会員数、ツアー公演数、初日・最終公演等を用いる。説明文からは、番手、総名義数、ランダム配布、同行条件、制作開放席、見切れ・注釈付き席、本人確認およびQRコード単位の条件等を構造化する。これらを表形式の機械学習モデルへ入力し、構造化情報のみを用いる場合と、説明文および市場固有の条件を追加する場合を比較する。評価では平均絶対誤差を主要指標とし、二乗平均平方根誤差、百分率誤差、決定係数、一定誤差範囲内の割合および価格帯別の偏りも確認する。

第二の目的は、各観測時点から1日、3日および7日以内の需要状態を確率として推定することである。需要状態は、掲載が継続する状態、売却を確認できた状態、削除された状態の三つに分ける。削除には、売却以外に出品者による取り下げや内容変更が含まれ得るため、売却と同一の需要として扱わない。モデルの良否は単純な正解率だけではなく、予測確率と実際の発生頻度が一致しているかを、log loss、Brier scoreおよび確率校正指標等によって評価する。これにより、現在の価格と条件を持つチケットが短期間に市場からどのような形で変化するかを定量化する。

第三の目的は、現在の出品と比較可能で、かつ一定額以上安い新規出品が1日、3日および7日以内に現れる確率を推定することである。比較可能性は、同一公演、枚数、券種、名義種別等の条件から定義する。単に市場全体の最低価格が下がるかを予測するのではなく、購入者が実際に代替候補として検討できる条件をそろえる。さらに、妥当価格モデル、需要状態モデルおよび安価な代替出品モデルの出力を統合し、「現在購入する」または「待つ」という判断方針を作成する。判断方針は、買い逃しを重視する場合、両者を均衡させる場合、節約を重視する場合に分け、履歴上の機会損失と行動件数から評価する。

第四の目的は、実際の利用時点では知り得ない情報を排除した検証方法を確立することである。価格や状態の予測では、予測時点より後の売却結果、市場価格および出品情報を特徴量へ含めない。同一またはほぼ同一の説明文を持つ出品が学習用と検証用に分かれて性能が過大評価されることを防ぎ、時間順に学習・検証を行う。また、モデル構成を決定した後に収集された新しいスナップショットを最終評価へ用い、市場時期や公演構成が変化した場合にも性能が維持されるかを確認する。データ取得の欠測、識別子変更および観測不能な売却時刻についても品質検査を行い、信頼できない期間は学習対象から除外する。

以上から、本研究では、どの情報がチケット価格の推定に有効か、説明文に含まれる市場固有条件が構造化情報のみのモデルを改善するか、現在の出品状態と将来の安価な代替出品をどの程度確率的に予測できるか、さらにそれらを統合することで購入時期の判断をどの程度支援できるかを研究課題とする。最終的には、経験や勘だけに依存していた二次流通市場の価格評価と購入判断を、データの根拠、予測誤差および不確実性を伴う再現可能な判断へ近づけることを目指す。ただし、モデルの出力は取引を保証または推奨するものではなく、市場分析上の参考情報として位置付ける。

\end{CJK}
\end{document}
